% Minimal preamble for Open Logic Project sources (pinned commit 1e960beff9ed7835bf3e3f1335e21af3439cd107).
%
% OLP's own master preamble cannot run on the project host: TeX Live 2025 there lacks
% prooftree.sty, ntheorem.sty, extarrows.sty, mathdots.sty, isomath.sty and subfiles.sty.
% Following the project's own rule ("if the author preamble needs packages the toolchain lacks,
% do not invent macros -- use a minimal working preamble"), this file declares, verbatim in
% spirit and numbering behaviour, exactly the theorem-like environments OLP defines in
% sty/open-logic-book-envs.sty (thm/lem/prop/cor/defn/ex/prob/rem/note/case/axiom/conv),
% plus the shorthand accents OLP uses (`!!a`, `!!b`, `!!{...}`).
%
% Author text is untouched; only the packaging is ours.
\documentclass[10pt,a4paper]{article}

\usepackage{amsmath}
\usepackage{amssymb}
\usepackage{amsthm}
\usepackage{graphicx}
\usepackage{hyperref}
\usepackage{enumitem}
\usepackage[framemethod=default]{mdframed}

% ---- OLP shorthand accents (from sty/open-logic.sty) ----
% NB: \! is taken by LaTeX (negative thin space), so we only define the !!a / !!b families
% and the \olup helper used by them.
\newcommand*{\olup}[1]{\MakeUppercase{#1}}
\def\!!!a#1{\olup{#1}}
\def\!a#1{\olup{#1}}
\def\!b#1{\olup{#1}}

% ---- theorem-like environments, same family numbering as OLP ----
\theoremstyle{plain}
\newtheorem{thm}{Theorem}[section]
\newtheorem{lem}[thm]{Lemma}
\newtheorem{prop}[thm]{Proposition}
\newtheorem{cor}[thm]{Corollary}
\newtheorem{axiom}{Axiom}

\theoremstyle{definition}
\newtheorem{defn}[thm]{Definition}
\newtheorem{ex}[thm]{Example}
\newtheorem{prob}{Problem}[section]
\newtheorem{probd}{Problem}
\newtheorem{conv}[thm]{Convention}

\theoremstyle{remark}
\newtheorem{rem}[thm]{Remark}
\newtheorem{note}[thm]{Note}
\newtheorem{case}[thm]{Case}

% ---- environments OLP defines as wrappers ----
\newenvironment{probtag}[1]{\begin{prob}}{\end{prob}}
\newenvironment{explain}{}{}
\newenvironment{explain*}{}{}
\newenvironment{editorial}{}{}
\newenvironment{digress}{}{}
\newenvironment{digress*}{}{}
\newenvironment{history}{}{}
\newenvironment{reference}{}{}
\newcommand*{\olfileid}[3]{}
\newcommand*{\olsection}[1]{\section*{#1}}
\newcommand*{\olchapter}[3]{}
\newcommand*{\olpart}[2]{}
\newcommand*{\olimport}[1]{}
\newcommand*{\olinclude}[1]{}
\newcommand*{\olpath}[1]{}
\newcommand*{\ollangid}[1]{}
\newcommand*{\problemtag}[1]{}
\newcommand*{\problemauthor}[1]{}
\newcommand*{\ollabel}[1]{\label{#1}}
% \olref takes 0-3 OPTIONAL arguments plus one mandatory label:
%     \olref{lab}  |  \olref[from]{lab}  |  \olref[from][to]{lab}  |  \olref[a][b][c]{lab}
% A single-argument definition silently swallows "[sth]" and emits a broken reference
% (measured on the host: the text came out as "??iso]x"). The chain below consumes the
% brackets and cross-references only the mandatory argument.
\makeatletter
\newcommand{\olref}{\@ifnextchar[{\olref@i}{\olref@final}}
\def\olref@i[#1]{\@ifnextchar[{\olref@ii}{\olref@final}}
\def\olref@ii[#1]{\@ifnextchar[{\olref@iii}{\olref@final}}
\def\olref@iii[#1]{\olref@final}
\def\olref@final#1{\ref{#1}}
\makeatother
\newcommand*{\olnewpage}{\clearpage}

% ===================== OLP compatibility layer (1) =====================

\makeatletter

% ---------------- proof trees (OLP's prooftree is a bussproofs derivative) ----------------
\providecommand{\AxiomC}[1]{}
\providecommand{\UnaryInfC}[1]{}
\providecommand{\BinaryInfC}[1]{}
\providecommand{\TrinaryInfC}[1]{}
\providecommand{\QuaternaryInfC}[1]{}
\providecommand{\QuinaryInfC}[1]{}
\providecommand{\RightLabel}[1]{}
\providecommand{\LeftLabel}[1]{}
\providecommand{\doubleLine}{}
\providecommand{\noLine}{}
\providecommand{\DisplayProof}{}
\providecommand{\fCenter}{}
\newenvironment{prooftree}{}{}
\newenvironment{ltprooftree}{}{}
\newenvironment{oltableau}{}{}

% ---------------- logical notation ----------------
\providecommand{\lif}{\mathbin{\rightarrow}}
\providecommand{\lforall}{\mathord{\forall}}
\providecommand{\lexists}{\mathord{\exists}}
\providecommand{\ltrue}{\top}
\providecommand{\lfalse}{\bot}
\providecommand{\Ltrue}{\top}
\providecommand{\Lfalse}{\bot}
\providecommand{\Sequent}{\mathrel{\Rightarrow}}
\providecommand{\Proves}{\mathrel{\vdash}}
\providecommand{\Entails}{\mathrel{\models}}
\providecommand{\Subst}[3]{#1[#2/#3]}
\providecommand{\sFmla}[2]{#1(#2)}
\providecommand{\pFmla}[2]{#1(#2)}
\providecommand{\formula}{\relax}
\providecommand{\True}{\top}
\providecommand{\False}{\bot}
\providecommand{\Args}[1]{#1}
\providecommand{\Ts}{\ensuremath{\top}}
\providecommand{\Fi}{\ensuremath{\bot}}
\providecommand{\lfalsum}{\bot}

% ---------------- objects, structures, metatheory ----------------
\providecommand{\Obj}{}
\providecommand{\Struct}{\ensuremath{\mathfrak{A}}}
\providecommand{\Domain}{\ensuremath{|\mathfrak{A}|}}
\providecommand{\Assign}{\ensuremath{s}}
\providecommand{\Sat}[2]{\ensuremath{#1 \models #2}}
\providecommand{\mSat}[3]{\ensuremath{#1 \models #2[#3]}}
\providecommand{\mModel}[2]{\ensuremath{#1 \models #2}}
\providecommand{\Lang}{\ensuremath{\mathcal{L}}}
\providecommand{\Th}{\ensuremath{\mathrm{Th}}}
\providecommand{\Log}{\ensuremath{\mathrm{L}}}
\providecommand{\tuple}[1]{\langle #1 \rangle}
\providecommand{\num}[1]{\underline{#1}}
\providecommand{\cfind}[1]{\varphi_{#1}}
\providecommand{\Nat}{\ensuremath{\mathbb{N}}}
\providecommand{\Mult}{\ensuremath{\mathrm{Mult}}}
\providecommand{\Add}{\ensuremath{\mathrm{Add}}}
\providecommand{\Exp}{\ensuremath{\mathrm{Exp}}}
\providecommand{\Zero}{\ensuremath{0}}
\providecommand{\Succ}{\ensuremath{\mathrm{S}}}
\providecommand{\lambd}{\ensuremath{\lambda}}
\providecommand{\bredpar}{\mathrel{\rightarrow_\beta}}
\providecommand{\bcd}[1]{\ensuremath{\ulcorner #1 \urcorner}}
\providecommand{\Set}[1]{\ensuremath{\{#1\}}}
\providecommand{\indcase}[1]{\par\noindent\emph{Case #1}.\quad}
\providecommand{\indfrm}[1]{\par\noindent\emph{Form #1}.\quad}
\providecommand{\iftag}[1]{#1}
\providecommand{\tagitem}[1]{\item[#1]}
\providecommand{\Elim}[1]{}
\providecommand{\Intro}[1]{}
\providecommand{\LeftR}[1]{}
\providecommand{\RightR}[1]{}
\providecommand{\TAss}[1]{}
\providecommand{\TRule}[1]{}

\makeatother
\providecommand{\olpath}{.}

% ===================== OLP compatibility layer (2) =====================
% --- second compatibility block: notation and shorthands used by the frozen OLP units ---
\providecommand{\CutCS}{\mathrm{Cut}}
\providecommand{\LeftSubproofLabel}{\relax}
\providecommand{\RightSubproofLabel}{\relax}
\providecommand{\maeh}{\relax}
\providecommand{\cheight}{\relax}
\providecommand{\ande}{\mathrm{and}_E}
\providecommand{\andi}{\mathrm{and}_I}
\providecommand{\ore}{\mathrm{or}_E}
\providecommand{\ori}{\mathrm{or}_I}
\providecommand{\typeof}{\mathord{:}}
\providecommand{\pValue}[1]{#1}
\providecommand{\Sent}{\mathrm{Sent}}
\providecommand{\cardplus}{\mathbin{+}}
\providecommand{\cardtimes}{\mathbin{\cdot}}
\providecommand{\cardexpo}{\mathbin{\hat{\ }}}
\providecommand{\cardnless}{\not<}
\providecommand{\cardle}{\leq}
\providecommand{\funfromto}[3]{#1\colon #2 \to #3}
\providecommand{\funimage}[1]{\operatorname{im}(#1)}
\providecommand{\fdefined}{\mathord{\downarrow}}
\providecommand{\canonord}{\operatorname{ord}}
\providecommand{\ordtimes}{\mathbin{\cdot}}
\providecommand{\ordeq}[2]{#1 = #2}
\providecommand{\ordneq}[2]{#1 \neq #2}
\providecommand{\ordtype}[1]{\operatorname{ord}(#1)}
\providecommand{\ordsucc}{\mathord{+1}}
\providecommand{\ordunion}{\cup}
\providecommand{\onesphere}{S^1}
\providecommand{\rotationsgroup}{\mathrm{SO}(2)}
\providecommand{\equivrep}[1]{[#1]}
\providecommand{\rlexless}{\prec}
\providecommand{\ZFC}{\mathrm{ZFC}}
\providecommand{\ZF}{\mathrm{ZF}}
\providecommand{\Rat}{\mathbb{Q}}
\providecommand{\Real}{\mathbb{R}}
\providecommand{\PosInt}{\mathbb{Z}^{+}}
\providecommand{\Id}{\mathrm{id}}
\providecommand{\Bin}{\mathrm{Bin}}
\providecommand{\pto}{\rightharpoonup}
\providecommand{\dom}{\operatorname{dom}}
\providecommand{\ran}{\operatorname{ran}}
\providecommand{\len}{\operatorname{len}}
\providecommand{\concat}{\mathbin{{}^\frown}}
\providecommand{\card}[1]{\lvert #1 \rvert}
\providecommand{\Pow}{\mathcal{P}}
\providecommand{\Setabs}[2]{\{\, #1 \mid #2 \,\}}
\providecommand{\comp}{\circ}
\providecommand{\Char}{\chi}
\providecommand{\Prf}{\operatorname{Prf}}
\providecommand{\Prov}{\operatorname{Prov}}
\providecommand{\Proj}{\operatorname{proj}}
\providecommand{\Gn}{\operatorname{gn}}
\providecommand{\gn}{\operatorname{gn}}
\providecommand{\Value}[1]{\operatorname{val}(#1)}
\providecommand{\Deduce}{\mathrm{D}}
\providecommand{\Dual}{\mathrm{d}}
\providecommand{\Ax}{\mathrm{Ax}}
\providecommand{\pSat}[3]{#1 \models #2\,[#3]}
\providecommand{\pAssign}{\sigma}
\providecommand{\bforall}{\forall}
\providecommand{\bexists}{\exists}
\providecommand{\liff}{\mathbin{\leftrightarrow}}
\providecommand{\FV}{\operatorname{FV}}
\providecommand{\UnaryInf}{\ensuremath{\mid}}
\providecommand{\BinaryInf}{\ensuremath{\mid}}
\providecommand{\depth}{\operatorname{depth}}
\providecommand{\oliflabeldef}[3]{#1}
\providecommand{\TMblank}{\sqcup}
\providecommand{\TMstroke}{\mid}
\providecommand{\TMstay}{\mathrm{S}}
\providecommand{\TMright}{\mathrm{R}}
\providecommand{\TMleft}{\mathrm{L}}
\providecommand{\TMendtape}{\triangleright}
\providecommand{\umin}{\mu}
\providecommand{\Sum}{\sum}
\providecommand{\citep}[1]{\cite{#1}}
\providecommand{\citet}[1]{\cite{#1}}
\providecommand{\citeauthor}[1]{\cite{#1}}
\providecommand{\citeyear}[1]{\cite{#1}}

% ===================== OLP compatibility layer (3) =====================
% --- third compatibility block: remaining notation seen in the frozen OLP units ---
\providecommand{\iso}{\cong}
\providecommand{\nszero}{\mathrm{0}}
\providecommand{\nssucc}{\mathrm{S}}
\providecommand{\nsplus}{\mathbin{+}}
\providecommand{\nssless}{<}
\providecommand{\nsless}{<}
\providecommand{\PL}{\mathrm{PL}}
\providecommand{\Nec}{\mathord{\Box}}
\providecommand{\RK}{\mathrm{RK}}
\providecommand{\ST}{\mathrm{ST}}
\providecommand{\SSubst}[3]{#1[#2/#3]}
\providecommand{\subst}[3]{#1[#2/#3]}
\providecommand{\tagfalse}{\mathrm{F}}
\providecommand{\cutr}{\mathrm{Cut}}
\providecommand{\cardneq}{\neq}
\providecommand{\issucc}[1]{\operatorname{succ}(#1)}
\providecommand{\tup}[1]{\langle #1 \rangle}
\providecommand{\Var}{\mathrm{Var}}
\providecommand{\Bew}{\operatorname{Bew}}
\providecommand{\Con}{\operatorname{Con}}
\providecommand{\Cons}{\operatorname{Cons}}
\providecommand{\pDeduce}{\mathrm{D}}
\providecommand{\psat}{\models}
\providecommand{\den}[1]{\llbracket #1 \rrbracket}
\providecommand{\Sem}[2]{\llbracket #1 \rrbracket_{#2}}
\providecommand{\DomainOf}[1]{\lvert #1 \rvert}
\providecommand{\Form}{\mathrm{Form}}
\providecommand{\Terms}{\mathrm{Term}}
\providecommand{\Atom}{\mathrm{Atom}}
\providecommand{\FVars}[1]{\operatorname{FV}(#1)}
\providecommand{\fv}[1]{\operatorname{fv}(#1)}
\providecommand{\tuplep}[1]{\langle #1 \rangle}
\providecommand{\restrict}[2]{#1 \upharpoonright #2}
\providecommand{\sto}{\rightarrow}
\providecommand{\FmlaSet}[1]{\mathrm{Fmla}_{#1}}
\providecommand{\Thm}{\mathrm{Thm}}
\providecommand{\Eval}{\operatorname{val}}
\providecommand{\blank}{\sqcup}
\providecommand{\blankTape}{\sqcup}
\providecommand{\Consis}{\operatorname{Con}}
\providecommand{\drs}{\mathrm{d}}
\providecommand{\VRule}{\mid}
\providecommand{\Bet}{\mathrm{B}}
\providecommand{\Seg}[1]{\mathrm{Seg}_{#1}}
\providecommand{\Enums}{\mathrm{En}}
\providecommand{\Z}{\mathbb{Z}}
\providecommand{\N}{\mathbb{N}}
\providecommand{\abs}[1]{\lvert #1 \rvert}
\providecommand{\floor}[1]{\lfloor #1 \rfloor}
\providecommand{\ceil}[1]{\lceil #1 \rceil}

% ===================== OLP compatibility layer (4) =====================
% --- fourth compatibility block: harvested from the per-candidate compile logs ---
\providecommand{\beredpar}{\mathrel{\rightarrow_\beta}}
\makeatletter
\def\bered{\@ifnextchar[{\bered@i}{\mathrel{\rightarrow_\beta}}}
\def\bered@i[#1]{\mathrel{\rightarrow_\beta}}
\def\bred{\@ifnextchar[{\bred@i}{\mathrel{\rightarrow_\beta}}}
\def\bred@i[#1]{\mathrel{\rightarrow_\beta}}
\makeatother
\providecommand{\beredone}{\mathrel{=_\beta}}
\providecommand{\bredone}{\mathrel{=_\beta}}
\providecommand{\xredone}{\mathrel{=_\beta}}
\providecommand{\aconvone}{\mathrel{=_\alpha}}
\providecommand{\cardless}{<}
\providecommand{\tagrefs}[1]{#1}
\providecommand{\ordplus}{\mathbin{+}}
\providecommand{\funrestrictionto}[2]{#1 \upharpoonright #2}
\providecommand{\fundefined}{\mathord{\uparrow}}
\providecommand{\Trm}{\mathrm{Trm}}
\providecommand{\OProv}{\operatorname{Prov}}
\providecommand{\OPrf}{\operatorname{Prf}}
\providecommand{\aeq}{\mathrel{=}}
\providecommand{\mClass}[1]{[#1]}
\providecommand{\equivclass}[1]{[#1]}
\providecommand{\pheight}{\operatorname{ht}}
\providecommand{\phy}{\operatorname{ht}}
\providecommand{\Cut}{\mathrm{Cut}}
\providecommand{\QuantRank}{\operatorname{qr}}
\providecommand{\elemequiv}{\mathrel{\equiv}}
\providecommand{\PIso}{\mathrm{PIso}}
\providecommand{\Theory}{\mathrm{Th}}
\providecommand{\supstrict}{\subsetneq}
\providecommand{\trcl}{\operatorname{trcl}}
\providecommand{\cref}[1]{\ref{#1}}
\providecommand{\Discharge}{\mathrm{Dis}}
\providecommand{\Weakening}{\mathrm{Wk}}
\providecommand{\ZFminus}{\mathrm{ZF}^{-}}
\makeatletter
\renewcommand{\Dual}{\@ifnextchar[{\Dual@i}{\ifmmode\mathrm{d}\else$\mathrm{d}$\fi}}
\def\Dual@i[#1]{\ifmmode\mathrm{d}[#1]\else$\mathrm{d}[#1]$\fi}
\makeatother
\providecommand{\RightR}{\mathrm{R}}
\providecommand{\LeftR}{\mathrm{L}}
\providecommand{\Intro}{\mathrm{I}}
\providecommand{\Axiom}{\mathrm{Ax}}
\providecommand{\formula}[1]{#1}
\providecommand{\stared}{\ast}
\providecommand{\oset}[2]{\{\, #1 \mid #2 \,\}}
\providecommand{\setof}[2]{\{\, #1 \mid #2 \,\}}
\providecommand{\pow}[1]{\mathcal{P}(#1)}
\providecommand{\ordpair}[2]{\langle #1, #2 \rangle}
\providecommand{\ordtriple}[3]{\langle #1, #2, #3 \rangle}
\providecommand{\union}{\cup}
\providecommand{\inter}{\cap}
\providecommand{\seq}[1]{\langle #1 \rangle}
\providecommand{\SuchThat}{\mid}
\providecommand{\where}{\mid}
\providecommand{\TVar}{\mathrm{TVar}}
\providecommand{\Bool}{\mathrm{Bool}}
\providecommand{\Val}{\operatorname{Val}}
\providecommand{\TL}[1]{\mathrm{TL}}
\providecommand{\xstared}{\ast}

% ===================== OLP compatibility layer (5) =====================
% --- fifth compatibility block: rescues the quarantined OLP candidates ---
\providecommand{\ident}{\operatorname{id}}
\providecommand{\setrank}{\operatorname{rank}}
\providecommand{\cardfont}{\mathfrak{c}}
\providecommand{\Complement}{\overline}
\providecommand{\QR}{\mathrm{QR}}
\providecommand{\becd}{\ensuremath{\ulcorner}}
\providecommand{\disjointsum}{\mathbin{\oplus}}
\providecommand{\cardeq}[2]{#1 = #2}
\providecommand{\xred}{\mathrel{\rightarrow}}
% --- sixth compatibility block: names reported by real compile logs ---
\providecommand{\ORefut}{\operatorname{Refut}}
\providecommand{\ORProv}{\operatorname{Prov}}
\providecommand{\cutrank}{\operatorname{cr}}
\providecommand{\Refut}{\operatorname{Refut}}
\providecommand{\DischargeRule}{\mathrm{Dis}}
\providecommand{\Contraction}{\mathrm{C}}
\providecommand{\bottomAlignProof}{\relax}
\providecommand{\FalseCl}{\bot}
\providecommand{\red}{\mathrel{\rightarrow}}
\providecommand{\varAssign}{\sigma}
\providecommand{\defiff}{\mathrel{\leftrightarrow}}
\providecommand{\FalseInt}{\bot}
\providecommand{\fn}{\operatorname{fn}}
\providecommand{\dotted}{\cdot}

% --- tag-macro semantic fix (2026-10-13, final): \iftag CHOOSES a branch and the OLP proof body is
% always the SECOND group (\iftag{tag}{body}, {tag}{body}{c}, {tag}{body}{c}{d}); a definition that
% prints the tag literal drops the proof body. Bodies contain blank lines, hence \long.
\makeatletter
\providecommand{\iftag}{}
\long\def\iftag#1#2{\@ifnextchar\bgroup{\iftag@three{#2}}{\iftag@stop{#2}}}
\long\def\iftag@three#1#2{\iftag@stop{#1}\@ifnextchar\bgroup{\iftag@eat}{\relax}}
\long\def\iftag@eat#1{}
\long\def\iftag@stop#1{#1}
\long\def\oliflabeldef#1#2{\@ifnextchar\bgroup{\@gobble}{\iftag@stop{#2}}}
\providecommand{\tagrefs}[1]{}
\providecommand{\settag}[1]{}
\makeatother
\begin{document}

% ==== OLP file: content/first-order-logic/axiomatic-deduction/axiomatic-deduction.tex ====
% Part: first-order-logic
% Chapter: axiomatic-deduction



\iftag{FOL}
      {\olchapter{fol}{axd}{Axiomatic \usetoken{P}{derivation}}}
      {\olchapter{pl}{axd}{Axiomatic \usetoken{P}{derivation}}}

\begin{editorial}
  No effort has been made yet to ensure that the material in this
  chapter respects various tags indicating which connectives and
  quantifiers are primitive or defined: all are assumed to be
  primitive, except $\liff$ which is assumed to be defined. If the FOL
  tag is true, we produce a version with quantifiers, otherwise
  without.
\end{editorial}

\olimport{rules-and-proofs}

\olimport{axioms-rules-propositional}

\iftag{FOL}{%
  \olimport{axioms-rules-quantifiers}
}{}

\olimport{proving-things}

\iftag{FOL}{%
  \olimport{proving-things-quant}
}{}

\olimport{proof-theoretic-notions}

\olimport{deduction-theorem}

\iftag{FOL}{%
  \olimport{deduction-theorem-quantifiers}
}{}

\olimport{provability-consistency}

\olimport{provability-propositional}

\iftag{FOL}{%
  \olimport{provability-quantifiers}
}{}

\olimport{soundness}

\iftag{FOL}{
  \olimport{identity}
}{}

\OLEndChapterHook


% ==== OLP file: content/first-order-logic/axiomatic-deduction/axioms-rules-propositional.tex ====
% Part: first-order-logic
% Chapter: axiomatic-deduction
% Section: axioms-rules-propositional



\iftag{FOL}
      {\olfileid{fol}{axd}{prp}}
      {\olfileid{pl}{axd}{prp}}

\olsection{Axioms and Rules for the Propositional Connectives}

\begin{defn}[Axioms]
The set of $\PAx$ of \emph{axioms} for the propositional connectives comprises
all \olup{formula}s of the following forms:
\begin{align}
  & (A \land B) \lif A \ollabel{ax:land1}\\
  & (A \land B) \lif B \ollabel{ax:land2}\\
  & A \lif (B \lif (A \land B)) \ollabel{ax:land3}\\
  & A \lif (A \lor B) \ollabel{ax:lor1}\\
  & A \lif (B \lor A) \ollabel{ax:lor2}\\
  & (A \lif C) \lif ((B \lif C) \lif ((A \lor B) \lif C)) \ollabel{ax:lor3}\\
  & A \lif (B \lif A) \ollabel{ax:lif1}\\
  & (A \lif (B \lif C)) \lif ((A \lif B) \lif (A \lif C)) \ollabel{ax:lif2}\\
  & (A \lif B) \lif ((A \lif \lnot B) \lif \lnot A) \ollabel{ax:lnot1}\\
  & \lnot A \lif (A \lif B) \ollabel{ax:lnot2}\\
  & \ltrue \ollabel{ax:ltrue}\\
  & \lfalse \lif A \ollabel{ax:lfalse1}\\
  & (A \lif \lfalse) \lif \lnot A \ollabel{ax:lfalse2}\\
  & \lnot\lnot A \lif A \ollabel{ax:dne}
\end{align}
\end{defn}

\begin{defn}[Modus ponens]
  If $B$ and $B \lif A$ already occur in \olup{derivation}, then $A$ is
  a correct inference step.
\end{defn}

We'll abbreviate the rule modus ponens as ``\MP.''


% ==== OLP file: content/first-order-logic/axiomatic-deduction/axioms-rules-quantifiers.tex ====
% Part: first-order-logic
% Chapter: axiomatic-deduction
% Section: axioms-rules-quantifiers



\olfileid{fol}{axd}{qua}

\olsection{Axioms and Rules for Quantifiers}

\begin{defn}[Axioms for quantifiers]
The \emph{axioms} governing quantifiers are
all instances of the following:
\begin{align}
\ollabel{ax:q1} & \lforall[x][B] \lif B(t), \\
\ollabel{ax:q2} & B(t) \lif \lexists[x][B].
\end{align}
for any closed term~$t$.
\end{defn}

\begin{defn}[Rules for quantifiers]
\item If $B \lif A(a)$ already occurs in the \olup{derivation} and $a$
  does not occur in~$\Gamma$ or~$B$, then $B \lif
  \lforall[x][A(x)]$ is a correct inference step.
\item If $A(a) \lif B$ already occurs in the \olup{derivation} and $a$
  does not occur in~$\Gamma$ or~$B$, then $\lexists[x][A(x)] \lif
  B$ is a correct inference step.
\end{defn}

We'll abbreviate either of these by ``\QR.''


% ==== OLP file: content/first-order-logic/axiomatic-deduction/deduction-theorem-quantifiers.tex ====
% Part: first-order-logic 
% Chapter: axiomatic-deduction 
% Section: deduction-theorem-quantifiers



\olfileid{fol}{axd}{ddq}

\olsection{The Deduction Theorem with Quantifiers}

\begin{thm}[Deduction Theorem]
\ollabel{thm:deduction-thm-q} If $\Gamma \cup \{A\} \Proves B$,
then $\Gamma \Proves A \lif B$.
\end{thm}

\begin{proof}
We again proceed by induction on the length
of the \olup{derivation} of $B$ from $\Gamma \cup \{A\}$.

The proof of the induction basis is identical to that in the proof
of~\olref[ded]{thm:deduction-thm}.

For the inductive step, suppose again that the \olup{derivation} of $B$
from $\Gamma \cup \{A\}$ ends with a step~$B$ which is justified by
an inference rule. If the inference rule is modus ponens, we proceed
as in the proof of \olref[ded]{thm:deduction-thm}. If the inference
rule is \QR, we know that $B \ident C \lif \lforall[x][D(x)]$ and
\olup{formula} of the form $C \lif D(a)$ appears earlier in the
\olup{derivation}, where $a$ does not occur in~$C$, $A$, or $\Gamma$. We
thus have that
\begin{align*}
  \Gamma \cup \{A\} & \Proves C \lif D(a),\\
  \intertext{and the induction hypothesis applies, i.e., we have that}
    \Gamma & \Proves A \lif (C \lif D(a)).\\
  \intertext{By}
  & \Proves (A \lif (C \lif D(a))) \lif ((A \land C) \lif D(a))\\
  \intertext{and modus ponens we get}
  \Gamma & \Proves (A \land C) \lif D(a).\\
  \intertext{Since the eigenvariable condition still applies, we can add a step to this \olup{derivation} justified by \QR, and get}
    \Gamma & \Proves (A \land C) \lif \lforall[x][D(x)].\\
    \intertext{We also have}
    & \Proves ((A \land C) \lif \lforall[x][D(x)]) \lif (A \lif (C \lif \lforall[x][D(x)]),\\
    \intertext{so by modus ponens,}
    \Gamma & \Proves A \lif (C \lif \lforall[x][D(x)]),
\end{align*}
i.e., $\Gamma \Proves A \lif B$.

We leave the case where $B$ is justified by the rule \QR, but is of
the form $\lexists[x][D(x)] \lif C$, as an exercise.
\end{proof}

\begin{prob}
  Complete the proof of \olref[fol][axd][ddq]{thm:deduction-thm-q}.
\end{prob}


% ==== OLP file: content/first-order-logic/axiomatic-deduction/deduction-theorem.tex ====
% Part: first-order-logic 
% Chapter: axiomatic-deduction 
% Section: deduction-theorem

% verification of properties of provability needed for maximally 
% consistent sets in the completeness chapter.



\iftag{FOL}
      {\olfileid{fol}{axd}{ded}}
      {\olfileid{pl}{axd}{ded}}

\olsection{The Deduction Theorem}

As we've seen, giving \olup{derivation}s in an axiomatic system is
cumbersome, and \olup{derivation}s may be hard to find. Rather than
actually write out long lists of \olup{formula}s, it is generally easier
to argue that such \olup{derivation}s exist, by making use of a few
simple results. We've already established three such results:
\olref[ptn]{prop:reflexivity} says we can always assert that $\Gamma
\Proves A$ when we know that $A \in
\Gamma$. \olref[ptn]{prop:monotonicity} says that if $\Gamma \Proves A$
then also $\Gamma \cup \{B\} \Proves A$. And
\olref[ptn]{prop:transitivity} implies that if $\Gamma \Proves A$ and
$A \Proves B$, then $\Gamma \Proves B$. Here's another simple
result, a ``meta''-version of modus ponens:

\begin{prop}
  \ollabel{prop:mp} If $\Gamma \Proves A$ and $\Gamma \Proves A \lif
  B$, then $\Gamma \Proves B$.
\end{prop}

\begin{proof}
  We have that $\{A, A \lif B\} \Proves B$:
  \begin{derivation}
    1. & $A$ & Hyp.\\
    2. & $A \lif B$ & Hyp.\\
    3. & $B$ & 1, 2, MP
  \end{derivation}
  By \olref[ptn]{prop:transitivity}, $\Gamma \Proves B$.
\end{proof}

The most important result we'll use in this context is the deduction
theorem:

\begin{thm}[Deduction Theorem]
\ollabel{thm:deduction-thm} $\Gamma \cup \{A\} \Proves B$ if and
only if $\Gamma \Proves A \lif B$.
\end{thm}

\begin{proof}
The ``if'' direction is immediate.  If $\Gamma \Proves A \lif B$
then also $\Gamma \cup \{A\} \Proves A \lif B$ by
\olref[ptn]{prop:monotonicity}.  Also, $\Gamma \cup \{A\} \Proves A$ by
\olref[ptn]{prop:reflexivity}. So, by \olref{prop:mp}, $\Gamma \cup
\{A\} \Proves B$.

For the ``only if'' direction, we proceed by induction on the length
of the \olup{derivation} of $B$ from $\Gamma \cup \{A\}$.

For the induction basis, we prove the claim for every \olup{derivation}
of length~$1$. !!^a{derivation} of~$B$ from $\Gamma \cup \{A\}$ of
length~$1$ consists of $B$ by itself; and if it is correct $B$ is
either $\in \Gamma \cup \{A\}$ or is an axiom.  If $B \in \Gamma$ or
is an axiom, then $\Gamma \Proves B$. We also have that $\Gamma
\Proves B \lif (A \lif B)$ by \olref[prp]{ax:lif1}, and
\olref{prop:mp} gives $\Gamma \Proves A \lif B$. If $B \in \{ A\}$
then $\Gamma \Proves A \lif B$ because the last \olup{sentence}~$A
\lif B$ is the same as $A \lif A$, and we have \olup{derive}d that in
\olref[pro]{ex:identity}.

For the inductive step, suppose \olup{derivation} of~$B$ from
$\Gamma \cup \{A\}$ ends with a step~$B$ which is justified by modus
ponens. (If it is not justified by modus ponens, $B \in \Gamma$, $B
\ident A$, or $B$ is an axiom, and the same reasoning as in the
induction basis applies.) Then some previous steps in the
\olup{derivation} are $C \lif B$ and $C$, for some \olup{formula}~$C$,
i.e., $\Gamma \cup \{A\} \Proves C \lif B$ and $\Gamma \cup \{A\}
\Proves C$, and the respective \olup{derivation}s are shorter, so the
inductive hypothesis applies to them. We thus have both:
\begin{align*}
  & \Gamma \Proves A \lif (C \lif B); \\
  & \Gamma \Proves A \lif C.
\end{align*}
But also
\[
\Gamma \Proves (A \lif (C \lif B)) \lif
((A\lif C)  \lif (A \lif B)),
\]
by \olref[prp]{ax:lif2}, and two applications of \olref{prop:mp} give
$\Gamma \Proves A \lif B$, as required.
\end{proof}

Notice how \olref[prp]{ax:lif1} and \olref[prp]{ax:lif2} were chosen
precisely so that the Deduction Theorem would hold.

The following are some useful facts about \olup{derivability}, which we
leave as exercises.

\begin{prop}
  \ollabel{prop:derivfacts}
\begin{enumerate}
\item $\Proves (A \lif B) \lif ((B \lif C)
  \lif (A \lif C)$; \ollabel{derivfacts:a}
\item If $\Gamma \cup \{ \lnot A\}
  \Proves \lnot B$ then $\Gamma \cup \{ B\} \Proves
  A$ (Contraposition); \ollabel{derivfacts:b}
\item  $\{ A, \lnot A\} \Proves
    B$ (Ex Falso Quodlibet, Explosion); \ollabel{derivfacts:c}
\item  $\{ \lnot\lnot A\} \Proves
  A$ (Double Negation Elimination);\ollabel{derivfacts:d}
\item If $\Gamma \Proves \lnot\lnot A$ then $\Gamma \Proves
  A$;\ollabel{derivfacts:e}
\end{enumerate}
\end{prop}

\tagprob{FOL}
\begin{prob}
  Prove \olref[fol][axd][ded]{prop:derivfacts}
\end{prob}
\tagendprob

\tagprob{notFOL}
\begin{prob}
  Prove \olref[pl][axd][ded]{prop:derivfacts}
\end{prob}
\tagendprob


% ==== OLP file: content/first-order-logic/axiomatic-deduction/identity.tex ====
% Part: first-order-logic
% Chapter: axiomatic-proofs
% Section: identity



\olfileid{fol}{axd}{ide}

\olsection{\usetoken{P}{derivation} with \usetoken{S}{identity}}

In order to accommodate $\eq$ in \olup{derivation}s, we simply add new
axiom schemas. The definition of \olup{derivation} and $\Proves$ remains
the same, we just also allow the new axioms.

\begin{defn}[Axioms for \olup{identity}]
\begin{align}
\ollabel{ax:id1} & \eq[t][t], \\
\ollabel{ax:id2} & \eq[t_1][t_2] \lif (B(t_1) \lif
  B(t_2)),
\end{align}
for any closed terms $t$, $t_1$, $t_2$.
\end{defn}

\begin{prop}\ollabel{prop:sound}
  The axioms \olref{ax:id1} and \olref{ax:id2} are valid.
\end{prop}

\begin{proof}
  Exercise.
\end{proof}

\begin{prob}
Prove \olref[fol][axd][ide]{prop:sound}.
\end{prob}

\begin{prop}\ollabel{prop:iden1}
 $\Gamma \Proves \eq[t][t]$, for any term $t$ and set~$\Gamma$.
\end{prop}

\begin{prop}\ollabel{prop:iden2}
  If $\Gamma \Proves A(t_1)$ and $\Gamma \Proves
  \eq[t_1][t_2]$, then $\Gamma \Proves A(t_2)$.
\end{prop}

\begin{proof}
The \olup{formula}
\[
(\eq[t_1][t_2] \lif (A(t_1) \lif A(t_2)))
\]
is an instance of~\olref{ax:id2}. The conclusion follows by two applications
of~\MP.
\end{proof}


% ==== OLP file: content/first-order-logic/axiomatic-deduction/proof-theoretic-notions.tex ====
% Part: first-order-logic
% Chapter: axiomatic-deduction
% Section: proof-theoretic-notions



\iftag{FOL}
      {\olfileid{fol}{axd}{ptn}}
      {\olfileid{pl}{axd}{ptn}}

\olsection{Proof-Theoretic Notions}

\begin{explain}
Just as we've defined a number of important semantic notions
(\iftag{FOL}{validity}{tautology}, entailment, satisfiability), we now
define corresponding \emph{proof-theoretic notions}.  These are not
defined by appeal to satisfaction of \olup{sentence}s in \olup{structure}s,
but by appeal to the \olup{derivability} or \olup{nonderivability} of
certain formulas.  It was an important discovery that these notions
coincide.  That they do is the content of the \emph{soundness} and
\emph{completeness theorems}.
\end{explain}

\begin{defn}[!!^{derivability}]
!!^a{formula}~$A$ is \emph{\olup{derivable}} from $\Gamma$, written
$\Gamma \Proves A$, if there is \olup{derivation} from~$\Gamma$ ending
in~$A$.
\end{defn}

\begin{defn}[Theorems]
!!^a{formula}~$A$ is a \emph{theorem} if there is \olup{derivation} of
$A$ from the empty set.  We write $\Proves A$ if $A$ is a theorem
and $\Proves/ A$ if it is not.
\end{defn}

\begin{defn}[Consistency]
A set $\Gamma$ of \olup{formula}s is \emph{consistent} if and only if
$\Gamma\Proves/ \lfalse$; it is \emph{inconsistent} otherwise.
\end{defn}

\begin{prop}[Reflexivity]
\ollabel{prop:reflexivity}
If $A \in \Gamma$, then $\Gamma \Proves A$.
\end{prop}

\begin{proof}
  The \olup{formula}~$A$ by itself is \olup{derivation} of~$A$ from~$\Gamma$.
\end{proof}

\begin{prop}[Monotonicity]
\ollabel{prop:monotonicity}
If $\Gamma \subseteq \Delta$ and $\Gamma \Proves A$, then $\Delta
\Proves A$.
\end{prop}

\begin{proof}
Any \olup{derivation} of $A$ from $\Gamma$ is also \olup{derivation} of
$A$ from~$\Delta$.
\end{proof}

\begin{prop}[Transitivity]
\ollabel{prop:transitivity}
If $\Gamma \Proves A$ and $\{A\} \cup \Delta \Proves
B$, then $\Gamma \cup \Delta \Proves B$.
\end{prop}

\begin{proof}
  Suppose $\{A\} \cup \Delta \Proves B$. Then there is
  \olup{derivation} $B_1$, \dots, $B_l = B$ from~$\{A\} \cup
  \Delta$. Some of the steps in that \olup{derivation} will be correct
  because of a rule which refers to a prior line~$B_i = A$. By
  hypothesis, there is \olup{derivation} of~$A$ from~$\Gamma$, i.e.,
  \olup{derivation}~$A_1$, \dots, $A_k = A$ where every $A_i$ is an
  axiom, \olup{element} of~$\Gamma$, or correct by a rule of
  inference. Now consider the sequence
  \[
  A_1, \dots, A_k = A, B_1, \dots, B_l = B.
  \]
  This is a correct \olup{derivation} of~$B$ from $\Gamma \cup \Delta$
  since every $B_i = A$ is now justified by the same rule which
  justifies~$A_k = A$.
\end{proof}

Note that this means that in particular if $\Gamma \Proves A$ and $A
\Proves B$, then $\Gamma \Proves B$. It follows also that if $A_1,
\dots, A_n \Proves B$ and $\Gamma \Proves A_i$ for each~$i$, then
$\Gamma \Proves B$.

\begin{prop}
\ollabel{prop:incons}
$\Gamma$ is inconsistent iff $\Gamma \Proves A$ for every~$A$.
\end{prop}

\begin{proof}
Exercise.
\end{proof}

\tagprob{FOL}
\begin{prob}
Prove \olref[fol][axd][ptn]{prop:incons}.
\end{prob}
\tagendprob

\tagprob{notFOL}
\begin{prob}
Prove \olref[pl][axd][ptn]{prop:incons}.
\end{prob}
\tagendprob

\begin{prop}[Compactness]
\ollabel{prop:proves-compact}
  \begin{enumerate}
  \item If $\Gamma \Proves A$ then there is a finite subset $\Gamma_0
    \subseteq \Gamma$ such that $\Gamma_0 \Proves A$.
  \item If every finite subset of~$\Gamma$ is
    consistent, then $\Gamma$ is consistent.
  \end{enumerate}
\end{prop}

\begin{proof}
  \begin{enumerate}
    \item If $\Gamma \Proves A$, then there is a finite sequence of
      \olup{formula}s $A_1$, \dots,~$A_n$ so that $A \ident A_n$ and
      each $A_i$ is either a logical axiom, \olup{element} of~$\Gamma$
      or follows from previous \olup{formula}s by modus ponens.  Take
      $\Gamma_0$ to be those $A_i$ which are in~$\Gamma$.  Then the
      \olup{derivation} is likewise \olup{derivation} from~$\Gamma_0$, and
      so $\Gamma_0 \Proves A$.
    \item This is the contrapositive of~(1) for the special case $A
      \ident \lfalse$.
\end{enumerate}
\end{proof}


% ==== OLP file: content/first-order-logic/axiomatic-deduction/provability-consistency.tex ====
% Part: first-order-logic
% Chapter: axiomatic-deduction
% Section: provability-consistency



\iftag{FOL}
      {\olfileid{fol}{axd}{prv}}
      {\olfileid{pl}{axd}{prv}}

\olsection{\usetoken{S}{derivability} and Consistency}

We will now establish a number of properties of the \olup{derivability}
relation.  They are independently interesting, but each will play a
role in the proof of the completeness theorem.

\begin{prop}\ollabel{prop:provability-contr}
  If $\Gamma \Proves A$ and $\Gamma \cup \{A\}$ is inconsistent,
  then $\Gamma$ is inconsistent.
\end{prop}

\begin{proof}
  If $\Gamma \cup \{A\}$ is inconsistent, then $\Gamma \cup \{A\}
  \Proves \lfalse$. By \olref[ptn]{prop:reflexivity}, $\Gamma \Proves
  B$ for every $B \in \Gamma$. Since also $\Gamma \Proves A$ by
  hypothesis, $\Gamma \Proves B$ for every $B \in \Gamma \cup
  \{A\}$. By \olref[ptn]{prop:transitivity}, $\Gamma \Proves
  \lfalse$, i.e., $\Gamma$ is inconsistent.
\end{proof}

\begin{prop}
\ollabel{prop:prov-incons}
$\Gamma \Proves A$ iff $\Gamma \cup \{\lnot A\}$ is inconsistent.
\end{prop}

\begin{proof}
First suppose $\Gamma \Proves A$. Then $\Gamma \cup \{\lnot A\}
\Proves A$ by \olref[ptn]{prop:monotonicity}. $\Gamma \cup \{\lnot
A\} \Proves \lnot A$ by \olref[ptn]{prop:reflexivity}.  We also have
$\Proves \lnot A \lif (A \lif \lfalse)$ by \olref[prp]{ax:lnot2}. So
by two applications of \olref[ded]{prop:mp}, we have $\Gamma \cup \{\lnot
A\} \Proves \lfalse$.

Now assume $\Gamma \cup \{\lnot A\}$ is inconsistent, i.e., $\Gamma
\cup \{\lnot A\} \Proves \lfalse$. By the deduction theorem, $\Gamma
\Proves \lnot A \lif \lfalse$. $\Gamma \Proves (\lnot A \lif
\lfalse) \lif \lnot\lnot A$ by \olref[prp]{ax:lfalse2}, so $\Gamma
\Proves \lnot\lnot A$ by \olref[ded]{prop:mp}. Since $\Gamma \Proves
\lnot\lnot A \lif A$ (\olref[prp]{ax:dne}), we have $\Gamma
\Proves A$ by \olref[ded]{prop:mp} again.
\end{proof}

\begin{prob}
Prove that $\Gamma \Proves \lnot A$ iff $\Gamma \cup \{A\}$ is
inconsistent.
\end{prob}

\begin{prop}\ollabel{prop:explicit-inc}
  If $\Gamma \Proves A$ and $\lnot A \in \Gamma$, then $\Gamma$ is
  inconsistent.
\end{prop}

\begin{proof}
  $\Gamma \Proves \lnot A \lif (A \lif \lfalse)$ by \olref[prp]{ax:lnot2}.
  $\Gamma \Proves \lfalse$ by two applications of \olref[ded]{prop:mp}.
\end{proof}


\begin{prop}\ollabel{prop:provability-exhaustive}
  If $\Gamma \cup \{A\}$ and $\Gamma \cup \{\lnot A\}$ are both
  inconsistent, then $\Gamma$ is inconsistent.
\end{prop}

\begin{proof}
Exercise.
\end{proof}

\tagprob{FOL}
\begin{prob}
  Prove \olref[fol][axd][prv]{prop:provability-exhaustive}
\end{prob}
\tagendprob

\tagprob{notFOL}
\begin{prob}
  Prove \olref[pl][axd][prv]{prop:provability-exhaustive}
\end{prob}
\tagendprob



% ==== OLP file: content/first-order-logic/axiomatic-deduction/provability-propositional.tex ====
% Part: first-order-logic
% Chapter: axiomatic-deduction
% Section: provability-propositional



\iftag{FOL}
      {\olfileid{fol}{axd}{ppr}}
      {\olfileid{pl}{axd}{ppr}}

\olsection{\usetoken{S}{derivability} and the Propositional Connectives}

\begin{explain}
  We establish that the \olup{derivability} relation~$\Proves$ of
  axiomatic deduction is strong enough to establish some basic facts
  involving the propositional connectives, such as that $A \land B
  \Proves A$ and $A, A \lif B \Proves B$ (modus ponens). These
  facts are needed for the proof of the completeness theorem.
\end{explain}

\begin{prop}\ollabel{prop:provability-land}
  \begin{enumerate}
  \item \ollabel{prop:provability-land-left} Both $A \land B \Proves
    A$ and $A \land B \Proves B$
  \item \ollabel{prop:provability-land-right} $A, B \Proves A \land B$.
  \end{enumerate}
\end{prop}

\begin{proof}
  \begin{enumerate}
    \item From \olref[prp]{ax:land1} and \olref[prp]{ax:land1} by
      modus ponens.
  \item From \olref[prp]{ax:land3} by two applications of
    modus ponens.
  \end{enumerate}
\end{proof}


\begin{prop}\ollabel{prop:provability-lor}
  \begin{enumerate}
  \item $A \lor B, \lnot A, \lnot B$ is inconsistent.
  \item Both $A \Proves A \lor B$ and $B \Proves A \lor B$.
  \end{enumerate}
\end{prop}

\begin{proof}
  \begin{enumerate}
  \item From \olref[prp]{ax:lnot1} we get $\Proves \lnot A \lif (A
    \lif \lfalse)$ and $\Proves \lnot B \lif (B \lif \lfalse)$. So
    by the deduction theorem, we have $\{\lnot A\} \Proves A \lif
    \lfalse$ and $\{\lnot B\} \Proves B \lif \lfalse$. From
    \olref[prp]{ax:lor3} we get $\{\lnot A, \lnot B\} \Proves (A
    \lor B) \lif \lfalse$. By the deduction theorem, $\{A \lor B,
    \lnot A, \lnot B\} \Proves \lfalse$.
  \item From \olref[prp]{ax:lor1} and \olref[prp]{ax:lor2} by modus
    ponsens.
  \end{enumerate}
\end{proof}

\begin{prop}\ollabel{prop:provability-lif}
  \begin{enumerate}
  \item \ollabel{prop:provability-lif-left}  $A, A \lif B \Proves B$.
  \item \ollabel{prop:provability-lif-right}
    Both $\lnot A \Proves A \lif B$ and $B \Proves A \lif B$.
  \end{enumerate}
\end{prop}

\begin{proof}
  \begin{enumerate}
  \item We can \olup{derive}:
    \begin{derivation}
      1. & $A$ & \Hyp\\
      2. & $A \lif B$ & \Hyp\\
      3. & $B$ & 1, 2, \MP
    \end{derivation}
    
  \item By \olref[prp]{ax:lnot2} and \olref[prp]{ax:lif1} and the
    deduction theorem, respectively.
  \end{enumerate}
\end{proof}


% ==== OLP file: content/first-order-logic/axiomatic-deduction/provability-quantifiers.tex ====
% Part: first-order-logic
% Chapter: axiomatic-deduction
% Section: provability-quantifiers

% verification of properties of provability needed for maximally
% consistent sets in the completeness chapter.



\olfileid{fol}{axd}{qpr}

\olsection{\usetoken{S}{derivability} and the Quantifiers}

\begin{explain}
  The completeness theorem also requires that axiomatic deductions
  yield the facts about~$\Proves$ established in this section.
\end{explain}

\begin{thm}
\ollabel{thm:strong-generalization} If $c$ is \olup{constant} not occurring
in $\Gamma$ or $A(x)$ and $\Gamma \Proves A(c)$, then $\Gamma
\Proves \lforall[x][A(x)]$.
\end{thm}

\begin{proof}
By the deduction theorem, $\Gamma \Proves \ltrue \lif A(c)$. Since
$c$ does not occur in $\Gamma$ or~$\top$, we get $\Gamma \Proves
\ltrue \lif \lforall[x][A(x)]$. By the deduction theorem again, $\Gamma \Proves
\lforall[x][A(x)]$.
\end{proof}

\begin{prop}
\ollabel{prop:provability-quantifiers}
\begin{tagenumerate}{prvEx,prvAll}
\tagitem{prvEx}{$A(t) \Proves \lexists[x][A(x)]$.}{}

\tagitem{prvAll}{$\lforall[x][A(x)] \Proves A(t)$.}{}
\end{tagenumerate}
\end{prop}

\begin{proof}
\begin{tagenumerate}{prvEx,prvAll}
\tagitem{prvEx}{By \olref[qua]{ax:q2} and the deduction theorem.}{}
\tagitem{prvAll}{By \olref[qua]{ax:q1} and the deduction theorem.}{}
\end{tagenumerate}
\end{proof}


% ==== OLP file: content/first-order-logic/axiomatic-deduction/provability.tex ====
% Part: first-order-logic 
% Chapter: axiomatic-deduction 
% Section: provability

% verification of properties of provability needed for maximally 
% consistent sets in the completeness chapter.



\iftag{FOL}
      {\olfileid{fol}{axd}{prv}}
      {\olfileid{pl}{axd}{prv}}

\olsection{Properties of \usetoken{S}{derivability}}

\begin{prop}[Monotonicity] If $\Gamma \subseteq \Delta$ and $\Gamma \Proves A$
then also $\Delta \Proves A$ \end{prop}

\begin{proof} Any finite $\Gamma_0 \subseteq \Gamma$ is also a finite
subset of~$\Delta$, so \olup{derivation} of $\Gamma_0 \Proves A$ also shows
$\Delta \Proves A$. \end{proof}

\begin{prop}\ollabel{prop:provability} $\Gamma \Proves A$ if and only if
$\Gamma\cup \{\lnot A \}$ is inconsistent. \end{prop}

\begin{prop} \begin{enumerate}
\item \ollabel{prop:provability-contr} If $\Gamma \Proves A$ and $\Gamma
\cup \{ A\} \Proves \lfalse$, then $\Gamma$ is inconsistent.

\item \ollabel{prop:provability-lnot} If $\Gamma \cup \{A\}
\Proves \lfalse$, then $\Gamma \Proves \lnot A$.

\item \ollabel{prop:provability-exhaustive} If $\Gamma \cup \{A\}
\Proves \lfalse$ and $\Gamma \cup \{\lnot A\} \Proves
\lfalse$, then $\Gamma \Proves \lfalse$.

\tagitem{prvOr}{\ollabel{prop:provability-lor-left} If $\Gamma \cup \{A\}
\Proves \lfalse$ and $\Gamma \cup \{B\} \Proves
\lfalse$, then $\Gamma \cup \{A \lor B\} \Proves \lfalse$.}{}

\tagitem{prvOr}{\ollabel{prop:provability-lor-right} If $\Gamma
\Proves A$ or $\Gamma \Proves B$, then $\Gamma
\Proves A \lor B$.}{}

\tagitem{prvAnd}{\ollabel{prop:provability-land-left} If $\Gamma
\Proves A \land B$ then $\Gamma \Proves A$ and
$\Gamma \Proves B$.}{}

\tagitem{prvAnd}{\ollabel{prop:provability-land-right} If $\Gamma
\Proves A$ and $\Gamma \Proves B$, then $\Gamma
\Proves A \land B$.}{}

\tagitem{prvIf}{\ollabel{prop:provability-mp} If $\Gamma \Proves
A$ and $\Gamma \Proves A \lif B$, then $\Gamma
\Proves B$.}{}

\tagitem{prvIf}{\ollabel{prop:provability-lif} If $\Gamma \Proves
\lnot A$ or $\Gamma \Proves B$, then $\Gamma \Proves
A \lif B$.}{} \end{enumerate} \end{prop}

\begin{proof} 
\begin{enumerate}
\item Assume that $\Gamma_0 \cup \{A\} \Proves \lfalse$.

\begin{derivation} 
1. & $\Gamma_0 \cup \{A\} \Proves \lfalse$ & hyp. \\
2. & $\Gamma_1 \Proves A$ & hyp.\\ 
3. & $\Gamma_0 \cup \Gamma_1 \Proves \lfalse$ & cut \\ 
\end{derivation}

Since $\Gamma_0 \subseteq \Gamma$ and $\Gamma_1 \subseteq \Gamma$,
$\Gamma_0 \cup \Gamma_1 \subseteq \Gamma$, hence $\Gamma \Proves \lfalse$.

\item We assume that $\Gamma_0 \cup\{\lnot A\} \Proves \lfalse$.

\begin{derivation} 
1. & $\Gamma_0 \cup\{\lnot A\} \Proves \lfalse$ & hyp. \\ 
2. & $\Gamma_0 \Proves \lnot A \lif \lfalse$ & Deduction Theorem\\
3. & $\Gamma_0 \Proves (\lnot A \lif \lfalse) \lif A$ & Ax0\\ 
4. & $\Gamma_0 \Proves A$ & MP 2, 3 \\ 
\end{derivation}

\item We assume that $\Gamma \cup \{ A \} \Proves \lfalse$ and that
$\Gamma_1 \cup \{\lnot A \} \Proves \lfalse$.

\begin{derivation} 
1. & $\Gamma_0 \cup\{A\} \Proves \lfalse$ & hyp. \\
2. & $\Gamma_0 \Proves A \lif \lfalse$ & Deduction Theorem \\ 
3. & $\Gamma_1 \cup \{\lnot A \} \Proves \lfalse$ & hyp\\ 
4. & $\Gamma_0 \Proves (A \lif \lfalse) \lif \lnot A$ & Ax0\\ 
5. & $\Gamma_0 \Proves \lnot A$ & MP 2, 4\\ 
6. & $\Gamma_0 \cup \Gamma_1 \Proves \lfalse$ & Cut 3, 5 \\
\end{derivation}

Since $\Gamma_0 \subseteq \Gamma$ and $\Gamma_1 \subseteq \Gamma$,
$\Gamma_0 \cup \Gamma_1 \subseteq \Gamma$. Hence $\Gamma \Proves \lfalse$.

% prop:provability-lor-left 
\tagitem{defOr}{}{
\iftag{probOr}{Exercise.}{Exercise.}}

% prop:provability-lor-right 
\tagitem{defOr}{}{ \iftag{probOr}{Exercise.}{
We assume that $\Gamma_0 \Proves A$.

\begin{derivation} 
1. & $\Gamma_0 \Proves A$ & hyp \\ 
2. & $\Gamma_0 \Proves A \lif (A \lor B)$ & Ax0\\ 
3. & $\Gamma_0 \Proves A \lor B$ & MP 1, 2 \\
\end{derivation}

Therefore $\Gamma \Proves A \lor B$. The proof for when $\Gamma \Proves
B$ is similar.}}

% prop:provability-land-left 
\tagitem{defAnd}{}{ \iftag{probAnd}{Exercise}{
Assume that $\Gamma_0 \Proves A \land B$

\begin{derivation} 
1. & $\Gamma_0 \Proves A \land B$ & hyp \\ 
2. & $\Gamma_0 \Proves (A\land B) \lif A $ & Ax0\\ 
3. & $\Gamma_0 \Proves A \lor B$ & MP 1, 2 \\
\end{derivation}

Hence, $\Gamma \Proves A$. A similar \olup{derivation} starting shows that
$\Gamma \Proves B$.}}

% prop:provability-land-right
\tagitem{defAnd}{}{\iftag{probAnd}{Exercise.}{Exercise.}}

% prop:provability-mp 
\tagitem{defIf}{}{ \iftag{probIf}{Exercise.}{ We
assume both that $\Gamma_0 \Proves A$ and $\Gamma_0 \Proves A \lif B $.

\begin{derivation} 
1. & $\Gamma_0 \Proves A$ & hyp \\ 
2. & $\Gamma_0 \Proves A \lif B $ & hyp\\ 
3. & $\Gamma_0 \Proves B$ & MP 1, 2\\
\end{derivation}

Since $\Gamma_0 \cup \Gamma_1 \subseteq \Gamma$, this means that $\Gamma
\Proves B$.}}

% prop:provability-lif 
\tagitem{defIf}{}{\iftag{probIf}{Exercise.}{First
suppose $\Gamma \Proves \lnot A$.
    
\begin{derivation} 
1. & $\Gamma_0 \Proves \lnot A$ & hyp \\ 
2. & $\Gamma_0 \Proves \lnot A \lif (A \lif B) $ & Ax0\\ 
3. & $\Gamma_0 \Proves A \lif B$ & MP 1,
2\\ \end{derivation}

The proof for $\Gamma \Proves B$ is similar:

\begin{derivation} 
1. & $\Gamma_0 \Proves B$ & hyp \\ 
2. & $\Gamma_0 \Proves B \lif (A \lif B) $ & Ax0\\ 
3. & $\Gamma_0 \Proves A \lif B$ & MP 1, 2\\
\end{derivation}}}

\end{enumerate} 
\end{proof}

\begin{probtag}{probOr,probAnd,probIf} Complete the proof of
\olref[fol][axd][prv]{prop:provability}. 
\end{probtag}

\begin{thm}[Generalization] \ollabel{thm:Generalization} If $\Gamma \Proves
A$ and $x$ is not free in any \olup{formula} in $\Gamma$, then $\Gamma
\Proves \lforall[x][A]$. 
\end{thm}

\begin{proof} By induction on $\mathsf{Thm}_\Gamma$: if $A$ is an axiom,
so is $\lforall[x][A]$. If $A\in \Gamma$ then $x$ is not free in $A$, so
that $A \lif \lforall[x][A]$ is an axiom (\textbf{Ax3}), and $\Gamma
\Proves \lforall[x][A]$ by MP. Now suppose $A$ follows by MP because
$\Gamma \Proves B$ and $\Gamma \Proves B \lif A$. By induction
hypothesis, $\Gamma \Proves \lforall[x][B]$ and $\Gamma \Proves
\lforall[x][( B \lif A)]$. But by \textbf{Ax2} also \[ \Gamma \Proves
\lforall[x][( B \lif A)] \lif (\lforall[x][ B] \lif \lforall[x][A]), \]
and two applications of MP give $\Gamma \Proves \lforall[x][A]$ as
desired. \end{proof}

\begin{thm}\ollabel{thm:WeakGen} (\emph{Weak Generalization on Constants})
If $\Gamma \Proves A$ and $c$ is \olup{constant} not in occurring in
$\Gamma$, then there is a variable $x$ not in $A$ such that $\Gamma
\Proves \lforall[x][\Subst{A}{x}{c}]$, and the proof does not involve $c$.
\end{thm}

\begin{proof} Let $A_1,\ldots,A_n$ be a proof of $A$ from $\Gamma$, so
that $A=A_n$. Pick a variable $x$ not in $A_1,\ldots,A_n$, and consider
the new sequence: $\Subst{A_1}{x}{c},\ldots,\Subst{A_n}{x}{c}.$ Such a
sequence is a proof of $\Subst{A}{x}{c}$ from $\Gamma$. In fact, for each
$i=1,\ldots,n$: \begin{itemize} \item if $A$ is an axiom, then so is
$\Subst{A_i}{x}{c}$; \item if $A_i \in \Gamma$, then since $c$ is not in
$\Gamma$, we have $\Subst{A_i}{x}{c} = A_i \in \Gamma$; \item if $A_i$
is obtained from $A_j$ and $A_j \lif A_i$ then $\Subst{A_i}{x}{c}$
follows by MP from $\Subst{A_j}{x}{c}$ and $\Subst{(A_j \lif A_i)}{x}{c}
= \Subst{A_j}{x}{c} \lif \Subst{A_i}{x}{c}$. \end{itemize} It is clear
that the \olup{constant} $c$ no longer occurs in the new sequence. Now let
$\Gamma'$ comprise those \olup{formula}s from $\Gamma$ that appear in the
\olup{derivation} of $\Subst{A}{x}{c}$; then $x$ is not free in any
\olup{formula} in $\Gamma'$ and $\Gamma' \Proves \Subst{A}{x}{c}$. By
generalization (\olref{thm:Generalization}) $\Gamma'
\Proves \lforall[x][\Subst{A}{x}{c}]$, whence by monotonicity $\Gamma \Proves
\lforall[x][\Subst{A}{x}{c}]$, as desired. 
\end{proof}

\begin{explain}
Our goal is to replace the requirement in \olref{thm:WeakGen} that $x$
does not occur in $A$ (at all) by the weaker requirements that $x$ is
not free in $A$ and it is \olup{free for} $c$ in $A$. Clearly this can
be accomplished by a change of bound \olup{variable} --- so that is what
we set out to prove first.
\end{explain}

\begin{lem}
\ollabel{lem:free-for} 
If $x$ is free for $c$ in $A$ and $y$ is not free in $A$, then $x$
is \olup{free for} $y$ in $\Subst{A}{y}{c}$.
\end{lem}

\begin{proof} If $x$ is not \olup{free for} $y$ in $\Subst{A}{y}{x}$, then
some free occurrence of $y$ in $\Subst{A}{y}{c}$ falls within the scope of
a quantifier $\lforall[x]$. But $y$ is not free in $A$, by hypothesis, so
all such occurrences come from the substitution $\Subst{}{y}{c}$. So some
occurrence of $c$ falls within the scope of $\lforall[x]$ and $x$ is not
\olup{free for} $c$ in $A$. 
\end{proof}

\begin{lem}[Change of Bound Variable]
\ollabel{lem:ChangeBdVar} 
If $x$ and $y$ are not free in $A$ and they are both \olup{free for} $c$
in $A$, then $\Proves \lforall[x][\Subst{A}{x}{c}] \equiv
\lforall[y][\Subst{A}{y}{c}]$ (and the proof does not involve $c$).
\end{lem}

\begin{proof} 
From the hypotheses, $y$ is \olup{free for} $c$ in $A$ and $x$ is not
free in $A$, so by the previous \olref{lem:free-for}, $y$ is
\olup{free for} $x$ in $\Subst{A}{x}{c}$. It follows that
$\lforall[x][\Subst{A}{x}{c} \lif \Subst{A}{x}{c} \Subst{}{y}{x}]$
is an axiom (\textbf{Ax1}). Since $x$ is not free in $A$, we have
$\Subst{A}{x}{c} \Subst{}{y}{x} = \Subst{A}{y}{c}$, so that 
\[
\Proves \lforall[x][\Subst{A}{x}{c}] \lif \Subst{A}{y}{c}, 
\] 
and by the Deduction Theorem $\lforall[x][\Subst{A}{x}{c}] \Proves
\Subst{A}{y}{c}$. Since $y$ is not free in $A$, it also not free in
$\lforall[x][\Subst{A}{x}{c}]$, so that by Generalization
$\lforall[x][\Subst{A}{x}{c}] \Proves \lforall[y][\Subst{A}{y}{c}]$,
and the Deduction Theorem again gives $\Proves
\lforall[x][\Subst{A}{x}{c}] \lif \lforall[y][\Subst{A}{y}{c}]$.
The proof of $\Proves \lforall[y][\Subst{A}{y}{c}] \lif \lforall[x]
[\Subst{A}{x}{c}]$ is perfectly symmetric, so that the conclusion
follows by Taut.
\end{proof}

\begin{thm}[Strong Generalization on Constants]
\ollabel{thm:StrongGen} 
If $\Gamma \Proves A$, the \olup{constant} $c$ does not occur in
$\Gamma$, and $x$ is not free in $A$ but it is \olup{free for} $c$ in $A$,
then $\Gamma \Proves \lforall[x][\Subst{A}{x}{c}]$. 
\end{thm}

\begin{proof} 
Since $\Gamma \Proves A$ and $c$ is not in $\Gamma$, then by
Weak Generalization there is \olup{variable} $y$ not in $A$ such that
$\Gamma \Proves \lforall[y][\Subst{A}{y}{c}]$. Since $y$ is not in $A$
(at all), it is not free in $A$ and it is also \olup{free for} $c$ in $A$;
if moreover (by hypothesis) $x$ is not free in $A$ and \olup{free for} $c$ in
$A$, then the requirements for a change of bound \olup{variable} are met, so
$\Proves \lforall[x][\Subst{A}{x}{c}] \equiv
\lforall[y][\Subst{A}{y}{c}]$, whence $\Gamma \Proves
\lforall[x][\Subst{A}{x}{c}]$. 
\end{proof}

\begin{thm} 
\ollabel{thm:provability-quantifiers}
\begin{tagenumerate}{prvEx,prvAll} 
\tagitem{prvEx}{If $\Gamma \Proves A(t)$ then $\Gamma \Proves
  \lexists[x][A(x)]$.}{} 
\tagitem{prvAll}{If $\Gamma \Proves \lforall[x][A(x)]$ then $\Gamma
  \Proves A(t)$.}{}
\end{tagenumerate} 
\end{thm}

\begin{proof} 
\begin{tagenumerate}{prvEx,prvAll} 
%%Need axioms for exists
\tagitem{prvEx}{Exercise.}{}
\tagitem{prvAll} {Assume that  $\Gamma_0 \Proves \lforall[x][A(x)]$.

\begin{derivation} 
1. & $\Gamma_0 \Proves \lforall[x][A(x)]$ & hyp. \\
2. & $\Gamma_0 \Proves \lforall[x][A(x)] \lif \Subst{A}{t}{x}$ & Ax1 \\ 
3. & $\Gamma_0 \Proves \Subst{A}{t}{x}$ & MP 1, 2 \\ 
\end{derivation}}{}

\end{tagenumerate} 
\end{proof}


% ==== OLP file: content/first-order-logic/axiomatic-deduction/proving-things-quant.tex ====
% Part: first-order-logic 
% Chapter: axiomatic-deduction 
% Section: proving-things-quant



\olfileid{fol}{axd}{prq}

\olsection{\usetoken{P}{derivation} with Quantifiers}

\begin{ex}
Let us give \olup{derivation} of $(\lforall[x][A(x)] \land
\lforall[y][B(y)]) \lif \lforall[x][(A(x) \land B(x))]$.

First, note that
\begin{align*}
  (\lforall[x][A(x)] \land \lforall[y][B(y)]) & \lif \lforall[x][A(x)]\\
  \intertext{is an instance of \olref[prp]{ax:land1}, and}
  \lforall[x][A(x)] & \lif A(a) \\
  \intertext{of \olref[qua]{ax:q1}. So, by \olref[pro]{prop:chain}, we know that}
  (\lforall[x][A(x)] \land \lforall[y][B(y)]) & \lif A(a) \\
  \intertext{is \olup{derivable}. Likewise, since}
  (\lforall[x][A(x)] \land \lforall[y][B(y)]) & \lif \lforall[y][B(y)] \qquad\text{and}\\
    \lforall[y][B(y)] & \lif B(a)\\
    \intertext{are instances of \olref[prp]{ax:land2} and \olref[qua]{ax:q1}, respectively,}
    (\lforall[x][A(x)] \land \lforall[y][B(y)]) & \lif B(a)\\
    \intertext{is derivable by \olref[pro]{prop:chain}. Using an appropriate instance of \olref[prp]{ax:land3} and two applications of~\MP, we see that}
    (\lforall[x][A(x)] \land \lforall[y][B(y)]) & \lif (A(a) \land B(a))\\
    \intertext{is derivable. We can now apply \QR{} to obtain}
(\lforall[x][A(x)] \land \lforall[y][B(y)]) & \lif \lforall[x][(A(x) \land B(x))].
\end{align*}
\end{ex}


% ==== OLP file: content/first-order-logic/axiomatic-deduction/proving-things.tex ====
% Part: first-order-logic
% Chapter: axiomatic-deduction
% Section: proving-things



\iftag{FOL}
      {\olfileid{fol}{axd}{pro}}
      {\olfileid{pl}{axd}{pro}}

\olsection{Examples of \usetoken{P}{derivation}}


\begin{ex}
Suppose we want to prove $(\lnot D \lor E) \lif (D \lif
E)$. Clearly, this is not an instance of any of our axioms, so we
have to use the \MP{} rule to \olup{derive} it. Our only rule is~MP, which
given $A$ and $A \lif B$ allows us to justify~$B$. One
strategy would be to use \olref[prp]{ax:lor3} with $A$ being $\lnot
D$, $B$ being $E$, and $C$ being $D \lif E$, i.e., the instance
\[
(\lnot D \lif (D \lif E)) \lif ((E \lif (D \lif E)) \lif ((\lnot
D \lor E) \lif (D \lif E))).
\]
Why? Two applications of MP yield the last part, which is what we
want.  And we easily see that $\lnot D \lif (D \lif E)$ is an
instance of \olref[prp]{ax:lnot2}, and $E \lif (D \lif E)$ is an
instance of \olref[prp]{ax:lif1}. So our \olup{derivation} is:
\begin{derivation}
  1. &  $\lnot D \lif (D \lif E)$ & \olref[prp]{ax:lnot2} \\
  2. & $(\lnot D \lif (D \lif E)) \lif {}$\\
  &\qquad $((E \lif (D \lif E)) \lif ((\lnot D \lor E) \lif (D \lif E)))$ & \olref[prp]{ax:lor3}\\
  3. & $(E \lif (D \lif E)) \lif ((\lnot D \lor E) \lif (D \lif E))$ & 1, 2, \MP\\
  4. & $E \lif (D \lif E)$ & \olref[prp]{ax:lif1}\\
  5. & $(\lnot D \lor E) \lif (D \lif E)$ & 3, 4, \MP
\end{derivation}
\end{ex}

\begin{ex}\ollabel{ex:identity}
Let's try to find \olup{derivation} of $D \lif D$.  It is not an
instance of an axiom, so we have to use \MP{} to \olup{derive} it.
\olref[prp]{ax:lif1} is an axiom of the form~$A \lif B$ to which we
could apply~\MP. To be useful, of course, the $B$ which \MP{} would
justify as a correct step in this case would have to be~$D \lif D$,
since this is what we want to \olup{derive}. That means $A$ would also
have to be $D$, i.e., we might look at this instance of
\olref[prp]{ax:lif1}:
\[
D \lif (D \lif D)
\]
In order to apply \MP, we would also need to justify the corresponding
second premise, namely~$A$. But in our case, that would be~$D$, and
we won't be able to \olup{derive}~$D$ by itself. So we need a different
strategy.

The other axiom involving just~$\lif$ is \olref[prp]{ax:lif2}, i.e.,
\[
(A \lif (B \lif C)) \lif ((A \lif B) \lif (A \lif C))
\]
We could get to the last nested conditional by applying \MP{}
twice. Again, that would mean that we want an instance of
\olref[prp]{ax:lif2} where $A \lif C$ is $D \lif D$, the \olup{formula} we
are aiming for. Then of course, $A$ and $C$ are both~$D$. How
should we pick~$B$ so that both $A \lif (B \lif C)$ and $A \lif
B$, i.e., in our case $D \lif (B \lif D)$ and $D \lif B$, are
also \olup{derivable}? Well, the first of these is already an instance of
\olref[prp]{ax:lif1}, whatever we decide $B$ to be. And $D \lif B$ would
be another instance of \olref[prp]{ax:lif1} if $B$ were $(D \lif D)$.
So, our \olup{derivation} is:
\begin{derivation}
  1. & $D \lif ((D \lif D) \lif D)$ & \olref[prp]{ax:lif1}\\
  2. & $(D \lif ((D \lif D) \lif D)) \lif {}$\\
  & \qquad $((D \lif (D \lif D)) \lif (D \lif D))$ & \olref[prp]{ax:lif2}\\
  3. & $(D \lif (D \lif D)) \lif (D \lif D)$ & 1, 2, \MP\\
  4. & $D \lif (D \lif D)$ & \olref[prp]{ax:lif1}\\
  5. & $D \lif D$ & 3, 4, \MP
\end{derivation}
\end{ex}

\begin{ex}\ollabel{ex:chain}
Sometimes we want to show that there is \olup{derivation} of some
\olup{formula} from some other \olup{formula}s~$\Gamma$. For instance, let's
show that we can \olup{derive} $A \lif C$ from $\Gamma = \{A \lif B,
B \lif C\}$.
\begin{derivation}
  1. & $A \lif B$ & \Hyp\\
  2. & $B \lif C$ & \Hyp\\
  3. & $(B \lif C) \lif (A \lif (B \lif C))$ & \olref[prp]{ax:lif1} \\
  4. & $A \lif (B \lif C)$ & 2, 3, \MP\\
  5. & $(A \lif (B \lif C)) \lif {}$\\
  & \qquad $((A \lif B) \lif (A \lif C))$ & \olref[prp]{ax:lif2}\\
  6. & $((A \lif B) \lif (A \lif C))$ & 4, 5, \MP\\
  7. & $A \lif C$ & 1, 6, \MP
\end{derivation}
The lines labelled ``\Hyp'' (for ``hypothesis'') indicate that the
\olup{formula} on that line is \olup{element} of~$\Gamma$.
\end{ex}

\begin{prop}
  \ollabel{prop:chain} If $\Gamma \Proves A \lif B$ and $\Gamma
  \Proves B \lif C$, then $\Gamma \Proves A \lif C$
\end{prop}

\begin{proof}
  Suppose $\Gamma \Proves A \lif B$ and $\Gamma \Proves B \lif
  C$. Then there is \olup{derivation} of $A \lif B$ from~$\Gamma$;
  and \olup{derivation} of~$B \lif C$ from~$\Gamma$ as well. Combine
  these into a single \olup{derivation} by concatenating them. Now add
  lines 3--7 of the \olup{derivation} in the preceding example. This is
  \olup{derivation} of $A \lif C$---which is the last line of the new
  \olup{derivation}---from~$\Gamma$. Note that the justifications of
  lines 4 and~7 remain valid if the reference to line number~1 is
  replaced by reference to the last line of the \olup{derivation} of~$A
  \lif B$, and reference to line number~2 by reference to the last
  line of the \olup{derivation} of~$B \lif C$.
\end{proof}

\begin{prob} 
Show that the following hold by exhibiting \olup{derivation}s from the
axioms:
\begin{enumerate} 
\item $(A \land B) \lif (B \land A)$
\item $((A \land B) \lif C) \lif (A \lif (B \lif C))$
\item $\lnot(A \lor B) \lif \lnot A$
\end{enumerate} 
\end{prob}




% ==== OLP file: content/first-order-logic/axiomatic-deduction/rules-and-proofs.tex ====
% Part: first-order-logic 
% Chapter: axiomatic-deduction 
% Section: rules-and-proofs



\iftag{FOL}
      {\olfileid{fol}{axd}{rul}}
      {\olfileid{pl}{axd}{rul}}

\olsection{Rules and \usetoken{P}{derivation}}

\begin{explain}
  Axiomatic \olup{derivation}s are perhaps the simplest \olup{derivation} system for
  logic. !!^a{derivation} is just a sequence of \olup{formula}s.  To
  count as \olup{derivation}, every \olup{formula} in the sequence must
  either be an instance of an axiom, or must follow from one or more
  \olup{formula}s that precede it in the sequence by a rule of inference.
  !!^a{derivation} \olup{derive}s its last \olup{formula}.
\end{explain}

\begin{defn}[!!^{derivability}]
If $\Gamma$ is a set of \olup{formula}s of $\Lang L$ then a
\emph{\olup{derivation}} from~$\Gamma$ is a finite sequence $A_1$,
\dots,~$A_n$ of \olup{formula}s where for each $i \le n$ one of the
following holds:
\begin{enumerate}
\item $A_i \in \Gamma$; or
\item $A_i$ is an axiom; or
\item $A_i$ follows from some $A_j$ (and $A_k$) with $j < i$ (and
  $k < i$) by a rule of inference.
\end{enumerate}
\end{defn}

What counts as a correct \olup{derivation} depends on which inference
rules we allow (and of course what we take to be axioms).  And an
inference rule is an if-then statement that tells us that, under
certain conditions, a step~$A_i$ in \olup{derivation} is a correct
inference step.

\begin{defn}[Rule of inference]
A \emph{rule of inference} gives a sufficient condition for what
counts as a correct inference step in \olup{derivation} from~$\Gamma$.
\end{defn}

For instance, since any one-element sequence $A$ with $A \in \Gamma$
trivially counts as \olup{derivation}, the following might be a very
simple rule of inference:
\begin{quote}
  If $A \in \Gamma$, then $A$ is always a correct inference step in
  any \olup{derivation} from~$\Gamma$.
\end{quote}
Similarly, if $A$ is one of the axioms, then $A$ by itself is
\olup{derivation}, and so this is also a rule of inference:
\begin{quote}
  If $A$ is an axiom, then $A$ is a correct inference step.
\end{quote}
It gets more interesting if the rule of inference appeals to
\olup{formula}s that appear before the step considered. The following
rule is called \emph{modus ponens:}
\begin{quote}
  If $B \lif A$ and $B$ occur higher up in the \olup{derivation},
  then~$A$ is a correct inference step.
\end{quote}
If this is the only rule of inference, then our definition of
\olup{derivation} above amounts to this: $A_1$, \dots,~$A_n$ is
\olup{derivation} iff for each $i \le n$ one of the following holds:
\begin{enumerate}
\item $A_i \in \Gamma$; or
\item $A_i$ is an axiom; or
\item for some $j < i$, $A_j$ is $B \lif A_i$, and for some $k < i$,
  $A_k$ is~$B$.
\end{enumerate}
The last clause says that $A_i$ follows from~$A_j$ ($B \lif A_i$) and $A_k$
($B$) by modus ponens. If we can go from $1$ to~$n$, and
each time we find \olup{formula}~$A_i$ that is either in~$\Gamma$, an
axiom, or which a rule of inference tells us that it is a correct
inference step, then the entire sequence counts as a correct
\olup{derivation}.

\begin{defn}[!!^{derivability}]
A \olup{formula}~$A$ is \emph{\olup{derivable}} from $\Gamma$, written
$\Gamma \Proves A$, if there is \olup{derivation} from $\Gamma$ ending
in $A$.
\end{defn}

\begin{defn}[Theorems]
!!^a{formula}~$A$ is a \emph{theorem} if there is \olup{derivation}
of~$A$ from the empty set.  We write $\Proves A$ if $A$ is a
theorem and $\Proves/ A$ if it is not.
\end{defn}



% ==== OLP file: content/first-order-logic/axiomatic-deduction/soundness.tex ====
% Part: first-order-logic
% Chapter: axiomatic-deduction
% Section: soundness



\iftag{FOL}
      {\olfileid{fol}{axd}{sou}}
      {\olfileid{pl}{axd}{sou}}
      
\olsection{Soundness}

\begin{explain}
!!^a{derivation} system, such as axiomatic deduction, is \emph{sound}
if it cannot \olup{derive} things that do not actually hold.  Soundness is
thus a kind of guaranteed safety property for \olup{derivation} systems.
Depending on which proof theoretic property is in question, we would
like to know for instance, that
\begin{enumerate}
\item every \olup{derivable}~$A$ is valid;
\item if $A$ is \olup{derivable} from some others~$\Gamma$, it is also a
  consequence of them;
\item if a set of \olup{formula}s~$\Gamma$ is inconsistent, it is
  unsatisfiable.
\end{enumerate}
These are important properties of \olup{derivation} system.  If any of them do
not hold, the \olup{derivation} system is deficient---it would \olup{derive} too much.
Consequently, establishing the soundness of \olup{derivation} system is of the
utmost importance.
\end{explain}

\begin{prop}
  If $A$ is an axiom, then
  \iftag{FOL}
    {$\Sat{M}{A}[s]$ for each \olup{structure}~$\Struct{M}$ and assignment~$s$.}
    {$\pSat{v}{A}$ for each \olup{valuation}~$\pAssign{v}$.}
\end{prop}

\begin{proof}
\iftag{FOL}{We have to verify that all the axioms are valid. For
  instance, here is the case for \olref[qua]{ax:q1}: suppose $t$ is
  \olup{free for} $x$ in $A$, and assume
  $\Sat{M}{\lforall[x][A]}[s]$. Then by definition of satisfaction,
  for each $\varAssign{s'}{s}{x}$, also $\Sat{M}{A}[s']$, and in particular
  this holds when $s'(x) = \Value{t}{M}[s]$. By
  \olref[syn][ext]{prop:ext-formulas},
  $\Sat{M}{\Subst{A}{t}{x}}[s]$. This shows that
  $\Sat{M}{(\lforall[x][A] \lif \Subst{A}{t}{x})}[s]$.}{Do truth
  tables for each axiom to verify that they are tautologies.}
\end{proof}

\begin{thm}[Soundness]
\ollabel{thm:soundness}
If $\Gamma \Proves A$ then $\Gamma \Entails A$.
\end{thm}

\begin{proof}
By induction on the length of the \olup{derivation} of $A$ from
$\Gamma$. If there are no steps justified by inferences, then all
\olup{formula}s in the \olup{derivation} are either instances of axioms or are
in~$\Gamma$. By the previous proposition, all the axioms are
\iftag{FOL}{valid}{tautologies}, and hence if $A$ is an axiom then
$\Gamma \Entails A$. If $A \in \Gamma$, then trivially $\Gamma
\Entails A$.

If the last step of the \olup{derivation} of~$A$ is justified by modus
ponens, then there are \olup{formula}s $B$ and $B \lif A$ in the
\olup{derivation}, and the induction hypothesis applies to the part of
the \olup{derivation} ending in those \olup{formula}s (since they contain at
least one fewer step justified by an inference).  So, by induction
hypothesis, $\Gamma \Entails B$ and $\Gamma \Entails B \lif
A$. Then $\Gamma \Entails A$ by
\iftag{FOL}
      {\olref[syn][sem]{thm:sem-deduction}}
      {\olref[pl][syn][sem]{thm:sem-deduction}}.

\iftag{FOL}{Now suppose the last step is justified by
\QR. Then that step has the form $C \lif \lforall[x][B(x)]$ and
there is a preceding step $C \lif B(c)$ with $c$ not in $\Gamma$,
$C$, or $\lforall[x][B(x)]$. By induction hypothesis, $\Gamma
\Entails C \lif B(c)$. By
\olref[syn][sem]{thm:sem-deduction}, $\Gamma \cup \{C\} \Entails
B(c)$.

Consider some structure~$\Struct{M}$ such that $\Sat{M}{\Gamma \cup
  \{C\}}$.  We need to show that $\Sat{M}{\lforall[x][B(x)]}$. Since
$\lforall[x][B(x)]$ is \olup{sentence}, this means we have to show that
for every variable assignment~$s$, $\Sat{M}{B(x)}[s]$
(\olref[syn][ass]{prop:sat-quant}). Since $\Gamma \cup \{C\}$
consists entirely of sentences, $\Sat{M}{D}[s]$ for all $D \in
\Gamma$ by \olref[syn][sat]{defn:satisfaction}.  Let $\Struct{M'}$ be
like $\Struct{M}$ except that $\Assign{c}{M'} = s(x)$.  Since $c$ does
not occur in~$\Gamma$ or~$C$, $\Sat{M'}{\Gamma \cup \{C\}}$ by
\olref[syn][ext]{cor:extensionality-sent}. Since $\Gamma \cup \{C\}
\Entails B(c)$, $\Sat{M'}{B(c)}$.  Since $B(c)$ is \olup{sentence},
$\Sat{M'}{B(c)}[s]$ by
\olref[syn][ass]{prop:sentence-sat-true}. $\Sat{M'}{B(x)}[s]$ iff
$\Sat{M'}{B(c)}[s]$ by \olref[syn][ext]{prop:ext-formulas} (recall that
$B(c)$ is just $\Subst{B(x)}{c}{x}$). So,
$\Sat{M'}{B(x)}[s]$. Since $c$ does not occur in~$B(x)$, by
\olref[syn][ext]{prop:extensionality}, $\Sat{M}{B(x)}[s]$. But $s$
was an arbitrary variable assignment, so
$\Sat{M}{\lforall[x][B(x)]}$. Thus $\Gamma \cup \{C\} \Entails
\lforall[x][B(x)]$. By \olref[syn][sem]{thm:sem-deduction}, $\Gamma
\Entails C \lif \lforall[x][B(x)]$.

The case where $A$ is justified by \QR{} but is of the form
$\lexists[x][B(x)] \lif C$ is left as an exercise.}{}
\end{proof}

\tagprob{FOL}
\begin{prob}
  Complete the proof of \olref[fol][axd][sou]{thm:soundness}.
\end{prob}
\tagendprob

\begin{cor}
\ollabel{cor:weak-soundness}
If $\Proves A$, then $A$ is \iftag{FOL}{valid}{a tautology}.
\end{cor}

\begin{cor}
\ollabel{cor:consistency-soundness}
If $\Gamma$ is satisfiable, then it is consistent.
\end{cor}

\begin{proof}
We prove the contrapositive.  Suppose that $\Gamma$ is not consistent.
Then $\Gamma \Proves \lfalse$, i.e., there is \olup{derivation} of
$\lfalse$ from~$\Gamma$. By \olref{thm:soundness}, any
\iftag{FOL}{\olup{structure}~$\Struct{M}$}{\olup{valuation}~$\pAssign{v}$}
that satisfies $\Gamma$ must satisfy~$\lfalse$.  Since
\iftag{FOL}{$\Sat/{M}{\lfalse}$ for every
  \olup{structure}~$\Struct{M}$}{$\pSat/{v}{\lfalse}$ for every
  \olup{valuation}~$\pAssign{v}$}, no
\iftag{FOL}{$\Struct{M}$}{$\pAssign{v}$} can satisfy $\Gamma$, i.e.,
$\Gamma$ is not satisfiable.
\end{proof}



\end{document}
